\section*{Chapitre \ref{ch:s1:asia-specific-equity-strategies} --- Stratégies d'actions spécifiques à l'Asie}

\begin{solution}{exo:s1:asia-specific-equity-strategies:1}
$\ln 1.5/\ln 1.1 = 4.25$, donc cinq jours sous une limite de variation de 10~\%, dont quatre bloqués~; $\ln 1.5/\ln 1.2 = 2.22$, donc trois jours sous une limite de 20~\%, dont deux bloqués.
\end{solution}

\begin{solution}{exo:s1:asia-specific-equity-strategies:2}
Le retard de la file vaut $\ln 1.15 - \ln 1.10 = 0.044$ en log, donc la probabilité d'exécution est $e^{-0.044/0.02} = 0.108$~: environ un ordre sur neuf est exécuté.
\end{solution}

\begin{solution}{exo:s1:asia-specific-equity-strategies:3}
$0.03 \times 0.6^4 = 0.0039$ subsiste, donc $0.03 - 0.0039 = 0.0261$, soit 2.6 points en log, s'est résorbé. Le chiffre de la semaine suivante dans la simulation, $-2.81~\%$,
ajoute la dérive et le bruit propres de la fair value.
\end{solution}

\begin{solution}{exo:s1:asia-specific-equity-strategies:4}
Un blocage poussé clôture au-dessus de la valeur de l'action~: les vendeurs sont heureux de vendre à la limite, si bien que chaque ordre en file est exécuté. L'ouverture suivante revient vers la
valeur, ce qui reprend une partie de la prime d'attention~: 1.67~\% contre 5.90~\% pour les blocages avec une file d'attente.
\end{solution}

\begin{solution}{exo:s1:asia-specific-equity-strategies:5}
Une limite plus large n'est atteinte que par des mouvements plus importants, et la part d'un grand mouvement que la limite retient est plus grande~: la file d'attente derrière chaque blocage croît avec
la largeur. La prime d'attention est la même pour toutes les largeurs.
\end{solution}

\begin{solution}{exo:s1:asia-specific-equity-strategies:6}
$0.0129 \times 1\,000 = 12.9$ par jour. La transaction exécutée rapporte $2.64~\% - 0.15~\% = 2.49~\%$ en moyenne, mais seulement sur les ordres exécutés, et ceux-ci s'exécutent surtout
dans les blocages les plus faibles.
\end{solution}

\begin{solution}{exo:s1:asia-specific-equity-strategies:7}
Les clôtures au plafond de variation tombent à 0.79~\% des jours-action~; le gain à l'ouverture suivante à 3.42~\% pour chaque clôture et 1.21~\% après exécution~; le retournement de la semaine à $-0.27~\%$. La
file d'attente seule vaut environ un point pour un acheteur en file~; le reste de la transaction, et presque tout le retournement, vient des acheteurs d'attention.
\end{solution}

\begin{solution}{exo:s1:asia-specific-equity-strategies:8}
Une clôture à la limite est un cours auquel presque personne qui voulait acheter n'a pu le faire~: la file ne s'exécute que dans les blocages les plus faibles, et non dans les plus forts. Pondéré
par les exécutions, le gain est de 2.64~\%, et l'attention idéalisée du modèle fait de ce chiffre même une borne supérieure. Un ratio de Sharpe suppose des exécutions issues de la file, des coûts et la
capacité de la file.
\end{solution}

\begin{solution}{pb:s1:asia-specific-equity-strategies:1}
\begin{enumerate}
\item La poursuite du mouvement d'une action bloquée dans le même sens à l'ouverture suivante~; des cours qui accélèrent vers la limite à mesure qu'ils s'en approchent.
\item Ils réduisent la variance observée, autocorrèlent positivement les rendements et concentrent des masses ponctuelles aux limites.
\item Vendre des actions le jour même de leur achat~; acheter à la limite et vendre le lendemain.
\item Rendements de fair value multipliés par 1.5~; clôtures à l'intérieur de la limite~; achats d'attention de 3~\% après les clôtures au plafond, décroissant de 0.6 par jour~; la moitié des cibles situées à moins de
2~\% de la limite poussées jusqu'au blocage~; exécutions en file qui baissent avec la file d'attente.
\item Achats d'attention par les particuliers, une hausse suivie d'un retournement sur la semaine, et 1.16~\% par jour pour les traders avisés qui achètent le jour de la limite et vendent le
lendemain.
\item Ils achètent les jours de clôture au plafond et vendent le lendemain~; leurs achats ce jour-là annoncent un retournement de long terme plus marqué~; leur poussée explique l'\omterm{def:s1:asia-specific-equity-strategies:magnet}{effet d'aimant}.
\item Les petits comptes ne savent pas prévoir les rendements et courent après les mouvements du jour~; les grands comptes prévoient correctement et opèrent à contre-courant, mais les coûts les absorbent.
\item Les \omterm{def:s1:asia-specific-equity-strategies:northbound}{flux Northbound} nets prédisent les rendements des actions de Shanghai connectées.
\item 11.24~\%, 1.29~\% et 0.10~\% des jours-action.
\item 4.51~\% et 2.64~\%.
\item La file s'exécute le plus lorsque le blocage est le plus faible, et le moins lorsque la file d'attente est la plus longue.
\item Il fait passer des jours de la bande 7.5~\%--9.5~\% à la limite~: les clôtures au plafond passent de 0.87~\% à 1.29~\% des jours-action.
\item \textbf{Résultat nommé.} À l'ouverture suivant une clôture au plafond, un acheteur à chaque clôture gagne 2.98~\%, 4.51~\% et 6.58~\% selon des limites de 5~\%, 10~\% et 20~\%
de variation de cours, un acheteur en file pondéré par les exécutions 2.12~\%, 2.64~\% et 3.20~\%, et le cours rend 1.76~\%, 2.81~\% et 5.27~\% sur la semaine suivante.
\item À 10~\%, la file seule rapporte 3.42~\% (1.21~\% exécutés)~; l'attention ajoute l'essentiel du reste et le reprend sur la semaine.
\item Elle suppose une vente à découvert, et l'emprunt de titres peut être rare ou interdit.
\item Des détentions trimestrielles publiées avec cinq jours de retard laissent un IC de 0.003 et un ratio de Sharpe de 0.26, contre 0.021 et 2.04 avec des données quotidiennes.
\item Exécutions issues de la file, ouverture comme sortie, droit de timbre et commissions, limites de variation en vigueur à chaque date, et seulement les actions empruntables pour les ventes à découvert.
\item Le signal des \omterm{def:s1:asia-specific-equity-strategies:northbound}{flux Northbound}, dont les données ont été tronquées par les bourses.
\item Les deux exploitent la \omterm{def:s1:flows:fit}{pression sur les prix} d'une demande prévisible~; ici la demande vient des achats d'attention et des ordres en file, et non des flux des fonds.
\item Elle vend du temps~: elle reporte un mouvement au lendemain, et ceux qui peuvent se mettre en file tôt ou vendre aux acheteurs du lendemain matin en tirent profit.
\end{enumerate}
\end{solution}

\begin{solution}{iq:s1:asia-specific-equity-strategies:1}
Les clôtures à la limite ne sont pas les cours auxquels l'action aurait été échangée~: les grands mouvements se répartissent sur plusieurs jours, si bien que les clôtures font apparaître une \omterm{def:s1:asia-specific-equity-strategies:continuation}{continuation après la limite de variation}
que personne n'aurait pu acheter à ces cours. Un signal de momentum ajusté sur elles trouve un rendement qu'il ne peut pas exploiter.
\end{solution}

\begin{solution}{iq:s1:asia-specific-equity-strategies:2}
Décider si la file d'attente et l'information justifient la mise en file~: rejoindre tard signifie ne s'exécuter que si le blocage se rompt, c'est-à-dire au moment où la transaction est la plus défavorable. Soit
se mettre en file tôt avec une taille que les exécutions peuvent supporter, soit attendre l'ouverture suivante et payer l'écart.
\end{solution}

\begin{solution}{iq:s1:asia-specific-equity-strategies:3}
Comparer la vitesse des mouvements de cours (ou le risque instantané d'atteindre le pas de cotation suivant) en fonction de la distance à la limite à celle des mêmes actions loin
de celle-ci, en contrôlant la volatilité~; l'\omterm{def:s1:asia-specific-equity-strategies:magnet}{effet d'aimant} se manifeste par une accélération près de la limite, plus forte à la limite haute.
\end{solution}

\begin{solution}{iq:s1:asia-specific-equity-strategies:4}
Exécutions à la limite uniquement depuis la file, selon la priorité~; aucune transaction au-delà des limites~; aucune vente le jour même d'actions achetées ce jour-là~; limites par compartiment et
par statut telles qu'elles étaient à chaque date~; droit de timbre sur les ventes.
\end{solution}

\begin{solution}{iq:s1:asia-specific-equity-strategies:5}
Une action bloquée à la baisse ne peut pas être vendue pendant autant de jours que la baisse l'exigerait~; la valeur en risque (VaR) calculée sur les rendements quotidiens la sous-estime~; un fonds dont les rachats sont
quotidiens peut devoir vendre ce qu'il peut, et non ce qu'il devrait.
\end{solution}

\begin{solution}{iq:s1:asia-specific-equity-strategies:6}
Avec $y = x - L$ suivant une loi exponentielle de moyenne $m$, $E[e^{-y/s}] = s/(s + m)$ et $E[y e^{-y/s}] = m s^2/(s + m)^2$, donc le gain d'un ordre exécuté est
$E[y e^{-y/s}]/E[e^{-y/s}] = m s/(m + s)$, inférieur à la fois à $m$ et à $s$~: les exécutions sélectionnent les petites files d'attente, et la moyenne inconditionnelle $m$ surestime le gain.
\end{solution}
